\chapter{Verification}
\label{ch:verify}

Three layers of tests, all run by \code{make test} from the nornRTOS directory
(\code{tools/test\_all.sh} in nornCore; about a minute with an ESP32 connected). Every
regression run includes the host tests, simavr for both AVR boards, the ESP32 in QEMU,
and the boards that are connected.

\section{Host unit tests}

The host port runs the kernel on \code{ucontext} with simulated ticks and simulated
interrupts, so a test drives every primitive through a scripted sequence of calls and
compares the trace of (task, event, value) with the expected one, deterministically.
They are built with \code{-Wconversion}, AddressSanitizer and UndefinedBehaviorSanitizer.

\begin{table}[htbp]
\centering\small
\begin{tabularx}{\linewidth}{@{}lX@{}}
\toprule
Test & What it pins down \\
\midrule
\code{test\_kernel}    & three tasks through the scheduler, events, delays, periodic waits, semaphores, mutexes \\
\code{test\_queue}     & queues and mailboxes: hand-over, overflow, timeouts, the ISR variants \\
\code{test\_notify}    & task notifications and software timers, one shot and periodic \\
\code{test\_pool}      & memory pools and the run time statistics \\
\code{test\_features}  & heap, event groups, suspend and resume, the tickless query \\
\code{test\_timeout}   & timeouts at the limit of 16 bit ticks: \code{RTOS\_MAX\_TIMEOUT} and beyond, \code{RTOS\_FOREVER - 1} \\
\code{test\_wrap}, \code{test\_wrap32} & the tick counter running through four wraps (16 and 32 bits) with every timed primitive in use, each interval checked to the tick \\
\code{test\_prio}      & the scheduler against mutex ceilings: an owner that blocked while holding, a waiter handed the mutex, nested ceilings \\
\code{test\_tie}       & the scheduling order at equal effective priority (section~\ref{sec:tie}) in the task-context pick and in the interrupt tail \\
\code{test\_lifecycle} & a task whose function returns is terminated and deaf; a reset slot starts a new task clean; a return while holding a mutex panics \\
\code{test\_robust}    & deadlock, stack canary, task watchdog, timer task, statistics, task info, trace, assertion; one child process per scenario \\
\code{test\_checks}    & every mandatory check with assertions compiled out, one child process per violation \\
\code{test\_api}       & the contract test: every public function called as the header documents it, return values and statistics checked \\
\code{test\_stress}    & six tasks pick random operations on every primitive with random timeouts while the idle hook injects random interrupts; the kernel invariants are checked after every operation; 8 seeds $\times$ 200\,000 ticks \\
\bottomrule
\end{tabularx}
\caption{The host unit tests in \code{nornCore/test/}.}
\end{table}

\code{make host-cov} builds the host tests with gcov and reports the line coverage per
kernel file: 98\,\% of the kernel, the rest being the reset paths of \code{rtos\_panic.c}
that the host cannot take and the stack realignment only the ESP32 needs. \code{make
analyze} runs GCC's static analyser over the kernel, the ports and the HAL with each
architecture's compiler and reports nothing.

\section{Test applications on the targets}

Each test case checks its own counters at the end and prints \code{RESULT tcNN PASS} or
\code{FAIL}, then halts, which ends a simulator run. tc01 to tc06 pass on a real Mega
2560 with the same counters as simavr, on the ESP32 board and in QEMU.

\begin{table}[htbp]
\centering\small
\begin{tabularx}{\linewidth}{@{}lX@{}}
\toprule
Case & What it runs \\
\midrule
tc01 & periodic task, events, a semaphore from an ISR, a ceiling mutex, the stack check \\
tc02 & a queue fed from an ISR with deliberate overflow, a mailbox between tasks \\
tc03 & notifications from an ISR and from timers, one shot and periodic software timers \\
tc04 & timing: context switch, interrupt latency, periodic jitter \\
tc05 & panic, reset and recovery: an application panic, a stack overflow, a late task, a starved idle task; the record checked after each reset \\
tc06 & floating point kept in registers across preemption; the lazy FPU context on the ESP32 \\
tc07 & ESP32: a second kernel instance on the other core answering requests over the shared memory link \\
tc08 & heap, event groups, suspend and resume, tickless idle together \\
tc09 & ESP32: code and constants running from flash through the cache \\
tc10 & the kernel benchmark in cycles (chapter~\ref{ch:perf}) \\
tc11 & stack need per task \\
tc12 & stack and interrupt entry latency against call depth \\
\bottomrule
\end{tabularx}
\caption{The test applications in \code{nornCore/applications/}.}
\end{table}

\section{Continuous integration and the simulators}

The host tests, simavr and QEMU run on every push. QEMU is not cycle accurate and stalls
when the host is busy: its timing figures mean nothing, and a heavy parallel build on
the host can fail the latency limits of tc04 and tc10. For stack use QEMU is exact. The
rule for every change to the kernel is: tests first, then the change, then
\code{tools/measure.sh} before and after, then the full regression.
